\section{Loss Functions}
\label{sec:losses}

\subsection{Motion Generation Objective}
\label{sec:vimo_losses}

\subsubsection{6D Rotation Representation}
Joint rotations are encoded as 6D vectors capturing the first two columns of the orthonormal rotation matrix. A 3D rotation $\mathbf{R} \in \mathbb{R}^{3 \times 3}$ is represented as $[\mathbf{a}_1, \mathbf{a}_2] \in \mathbb{R}^6$, where $\mathbf{a}_1$ and $\mathbf{a}_2$ are orthonormal column vectors. The third column is recovered via cross product $\mathbf{a}_3 = \mathbf{a}_1 \times \mathbf{a}_2$. This representation is continuous, avoids the gimbal lock of Euler angles, and eliminates the unit-norm constraint of quaternions, making it better suited for gradient-based optimization~\cite{6drot2020}.

\subsubsection{Standard Regularizations}
Three auxiliary losses enforce physical plausibility.

\textit{3D Positions Loss} penalizes discrepancy in global joint positions:
\[
\mathcal{L}_{\text{joints}} = \frac{1}{S}\sum_{i=1}^{S} \| \text{FK}(\mathbf{R}_0^i) - \text{FK}(\hat{\mathbf{R}}_0^i) \|_2^2,
\]
where $\text{FK}$ denotes forward kinematics converting rotations to joint positions.

\textit{Velocity Loss} enforces temporal smoothness via finite differences:
\[
\mathcal{L}_{\text{vel}} = \frac{1}{S-1}\sum_{i=1}^{S-1} \| (\mathbf{R}_0^{i+1} - \mathbf{R}_0^i) - (\hat{\mathbf{R}}_0^{i+1} - \hat{\mathbf{R}}_0^i) \|_2^2,
\]
approximating angular velocity in 6D space.

\textit{Foot Contact Loss} prevents foot sliding by masking joint velocities with predicted contact labels:
\[
\mathcal{L}_{\text{foot}} = \frac{1}{S-1}\sum_{i=1}^{S-1} \| (\text{FK}(\hat{\mathbf{R}}_0^{i+1}) - \text{FK}(\hat{\mathbf{R}}_0^i)) \cdot \hat{\mathbf{f}}_i \|_2^2,
\]
where $\hat{\mathbf{f}}_i \in \{0,1\}^4$ indicates left/right foot/toe ground contact. If $\hat{f}_i = 1$, the velocity must be near zero.

Traditional methods compute $\hat{\mathbf{f}}_i$ via heuristic thresholds on foot velocity. ViMo-Flow learns $\hat{\mathbf{f}}_i$ as part of the model output without direct supervision; the only signal comes from $\mathcal{L}_{\text{foot}}$, encouraging the model to predict contact labels that minimize sliding. This self-supervised approach avoids arbitrary thresholds and adapts to the model's own motion predictions.

\subsubsection{SNR Integration}
Training directly predicts clean motion $\hat{\mathbf{m}}_0$ rather than noise. The simple diffusion loss, together with all auxiliary losses ($\mathcal{L}_{\text{joints}}$, $\mathcal{L}_{\text{vel}}$, $\mathcal{L}_{\text{foot}}$), is reweighted by timestep-dependent SNR to balance contributions across the noise schedule. The SNR at timestep $t$ is $\text{SNR}_t = \bar{\alpha}_t / (1-\bar{\alpha}_t)$. Pure SNR weighting $w_t = \text{SNR}_t$ produces a steep cliff where most timesteps contribute near-zero gradient, yielding poor sample efficiency.

Our implementation adopts a min-cap strategy that automatically selects the appropriate weighting across the diffusion process:
\[
w_t = \min\left( \text{SNR}_t,\ \sqrt{1 - \frac{t}{T}} \right).
\]
This formulation provides a smooth, monotonic ceiling: at early timesteps (high noise), the square-root decay caps the potentially large SNR$_t$ to prevent gradient explosion; at late timesteps (low noise), SNR$_t$ becomes small and dominates, avoiding vanishing gradients. The min operation preserves the natural emphasis on coarse-structure recovery while keeping all timesteps active during training.

The overall training objective becomes:
\[
\mathcal{L} = \mathbb{E}_t [ w_t \|\mathbf{m}_0 - \hat{\mathbf{m}}_0\|_2^2 ] + \lambda_1 \mathbb{E}_t [ w_t \mathcal{L}_{\text{joints}} ] + \lambda_2 \mathbb{E}_t [ w_t \mathcal{L}_{\text{vel}} ] + \lambda_3 \mathbb{E}_t [ w_t \mathcal{L}_{\text{foot}} ].
\]
Weights $w_t$ are normalized per batch to stabilize training. This square-root-capped Min-SNR strategy preserves sample efficiency, prevents collapse to mean poses, and accelerates convergence by emphasizing timesteps where the signal structure is most informative~\cite{minsnr2024}.

\subsection{Adversarial Control Objectives}
\label{sec:physics_losses}

\subsubsection{Discriminator Training Objective}
Each body-group discriminator $D_k$ is trained with hinge loss to separate reference (real) from policy-generated (fake) motion windows. Since each discriminator is implemented as an ensemble of $N=32$ scalar heads, the training objective is written head-wise as
\begin{equation}
    \mathcal{L}_{D_k}
    =
    \frac{1}{N}\sum_{i=1}^{N}
    \Bigl(
        \underbrace{\mathbb{E}_{\mathbf{o} \sim \mathrm{ref}}\!\bigl[\max(0,\, 1 - D_k^{(i)}(\mathbf{o}))\bigr]}_{\text{real loss}}
        +
        \underbrace{\mathbb{E}_{\mathbf{o} \sim \pi}\!\bigl[\max(0,\, 1 + D_k^{(i)}(\mathbf{o}))\bigr]}_{\text{fake loss}}
    \Bigr)
    +
    \lambda_{\text{gp}} \cdot \mathrm{GP}_k,
    \label{eq:disc_loss}
\end{equation}
where the gradient penalty~\cite{wgangp2017}
\[
    \mathrm{GP}_k = \mathbb{E}_{\hat{\mathbf{o}}}\!\bigl[\bigl(\|\nabla_{\hat{\mathbf{o}}} \sum_{i=1}^{N} D_k^{(i)}(\hat{\mathbf{o}})\|_2 - 1\bigr)^2\bigr],
    \quad \hat{\mathbf{o}} = \alpha\,\mathbf{o}_{\text{ref}} + (1{-}\alpha)\,\mathbf{o}_\pi,\; \alpha \sim U(0,1),
\]
enforces Lipschitz smoothness with $\lambda_{\text{gp}} = 10$. Hinge loss bounds scores to $[-1,1]$: real samples are pushed toward $+1$, fake samples toward $-1$. Once a sample is correctly classified with margin $\geq 1$, the corresponding hinge term saturates, preventing discriminator over-training on already-separated samples and preserving informative rewards~\cite{iccgan2021}.

The reward signal per discriminator is the mean over 32 ensemble heads:
\[
    r_t^{(k)} = \frac{1}{32}\sum_{i=1}^{32} \mathrm{clip}\!\bigl(D_k^{(i)}(\mathbf{o}_{t-H_k:t}),\, -1,\, 1\bigr).
\]

\subsubsection{Policy Optimization via PPO}
The policy is optimised with PPO~\cite{ppo2017}. Advantage estimates $\hat{A}_t$ are computed via GAE:
\[
    \hat{A}_t^{(i)} = \sum_{l=0}^{H_{\text{steps}}-1-t}(\gamma\lambda)^l\, \delta_{t+l}^{(i)}, \qquad
    \delta_t^{(i)} = r_t^{(i)} + \gamma V^{(i)}(\mathbf{s}_{t+1}) - V^{(i)}(\mathbf{s}_t).
\]
For multi-objective training, each advantage stream is standardized independently across the rollout batch:
\[
    \bar{A}_t^{(i)} = \frac{\hat{A}_t^{(i)} - \mu_i}{\sigma_i + \varepsilon},
\]
where $(\mu_i,\sigma_i)$ are the batch mean and standard deviation of objective $i$. The weighted advantage is then
\[
    \bar{A}_t = \sum_i w_i \bar{A}_t^{(i)}.
\]
The clipped surrogate objective is
\[
    \mathcal{L}^{\text{clip}} = -\mathbb{E}_t\!\left[\min\!\left(\rho_t\,\bar{A}_t,\; \mathrm{clip}(\rho_t,\, 1{-}\epsilon,\, 1{+}\epsilon)\,\bar{A}_t\right)\right],
    \quad \rho_t = \frac{\pi_\theta(\mathbf{a}_t|\mathbf{s}_t)}{\pi_{\theta_{\text{old}}}(\mathbf{a}_t|\mathbf{s}_t)},
\]
with $\epsilon = 0.2$. The value loss is
\[
    \mathcal{L}^{\text{value}} = \sum_i \bigl(V_\theta^{(i)}(\mathbf{s}_t) - R_t^{(i)}\bigr)^2,
\]
where targets $R_t^{(i)}$ are normalized per component via DiagonalPopArt to handle the varying scales of discriminator and task rewards. This per-objective normalization followed by weighted combination matches the multi-critic implementation used in training.

\subsubsection{Auxiliary Regularisations}

\textbf{Bilateral symmetry loss.} An optional $\ell_2$ penalty between mirrored sagittal-plane joint pairs $\mathcal{P} = \{(i,j)\}$ (11 pairs: shoulders, elbows, hips, knees, ankles):
\[
    \mathcal{L}_{\text{sym}} = \frac{\lambda_{\text{sym}}}{|\mathcal{P}|}\sum_{(i,j)\in\mathcal{P}} \|\sigma_j\,\mu_i - \mu_j\|^2,
\]
where $\mu_i,\mu_j$ are policy mean outputs and $\sigma_j \in \{-1,+1\}$ accounts for lateral-axis sign flips~\cite{symlocom2018}. The full PPO objective is:
\[
    \mathcal{L} = \mathcal{L}^{\text{clip}} + c_v\,\mathcal{L}^{\text{value}} + \mathcal{L}_{\text{sym}}.
\]

\textbf{Phase-conditioned observations.} When enabled, the normalised cycle position
\[
    \phi_t = \frac{t \bmod T_{\text{cycle}}}{T_{\text{cycle}}} \in [0,1)
\]
is encoded as $(\sin 2\pi\phi_t,\, \cos 2\pi\phi_t)$ and appended to the policy observation, providing an explicit temporal anchor for periodic motions.

\textbf{Termination penalty.} Episodes terminate when root body height drops below $h_{\min}$ for $N_{\text{grace}}$ consecutive frames (fall detection). Terminated transitions receive a fixed negative reward $r_{\text{term}}$ to discourage unstable policies. Looped reference motions wrap cyclically, and the maximum number of cycles per episode is configurable.
